\documentclass[11pt]{article}

\usepackage[T1]{fontenc}
\usepackage[utf8]{inputenc}
\usepackage{lmodern}
\usepackage[a4paper,margin=30mm]{geometry}
\usepackage{amsmath,amssymb,amsthm,mathtools}
\usepackage{booktabs,longtable,array}
\usepackage{enumitem}
\usepackage{microtype}
\usepackage{pdflscape}
\usepackage{xcolor}
\usepackage{hyperref}
\usepackage[capitalise,noabbrev]{cleveref}

\numberwithin{equation}{section}

\newtheorem{theorem}{Theorem}[section]
\newtheorem{proposition}[theorem]{Proposition}
\newtheorem{lemma}[theorem]{Lemma}
\newtheorem{corollary}[theorem]{Corollary}
\theoremstyle{definition}
\newtheorem{definition}[theorem]{Definition}
\theoremstyle{remark}
\newtheorem{remark}[theorem]{Remark}

\newcommand{\dd}{\,\mathrm d}
\newcommand{\one}{\mathbf 1}
\newcommand{\PhiH}{\Phi_H}
\newcommand{\join}{\mathbin{\vee}}
\newcommand{\ip}[2]{\langle #1,#2\rangle}
\newcommand{\Atlas}{\operatorname{Atlas}}
\newcommand{\Tail}{\operatorname{Tail}}

\title{A Complete Six-Vertex Classification for a Chromatic Graphon Lower Bound}
\author{Author names to be confirmed\thanks{The author list and affiliations will be inserted after confirmation.}}
\date{Draft of \today}

\begin{document}
\maketitle

\begin{abstract}
For a finite graph $H$, let
\[
 \Phi_H(p)=(1-p)^{v(H)}\chi_H\!\left(\frac1{1-p}\right),
\]
where the value at $p=1$ is defined by polynomial continuation.  We classify
the inequality $t(H,W)\geq\Phi_H(p)$ on the interval
$p\geq1-1/(\chi(H)-1)$ for every non-bipartite graph $H$ on at most six
vertices having neither isolated vertices nor an isolated-edge component.
There are $117$ such isomorphism classes.  The inequality holds for $94$ and
fails for $23$.  The positive side combines odd-cycle and pure-chordal
theorems, rooted and link inequalities, closure operations, row-specific
analytic arguments, and exact rational interval sum-of-squares certificates.
Every negative row has an explicit rational step-graphon witness; nineteen
come from tensor products of balanced Tur\'{a}n graphons.  Lean sources state
the row theorem for arbitrary probability spaces and arbitrary measurable
symmetric $[0,1]$-valued kernels.  Exact scripts regenerate the catalogue and
the counterexample arithmetic.
\end{abstract}

\section{Introduction}

Let $W\colon\Omega^2\to[0,1]$ be a symmetric measurable kernel on a
probability space $(\Omega,\mu)$, in the graphon framework of
Lov\'asz~\cite{Lovasz2012}.  For a finite simple graph $H$, its
homomorphism density is
\begin{equation}
 t(H,W)=\int_{\Omega^{V(H)}}
 \prod_{uv\in E(H)}W(x_u,x_v)
 \prod_{u\in V(H)}\dd\mu(x_u).
 \label{eq:hom-density}
\end{equation}
The edge density is $p=t(K_2,W)$.  At a balanced complete $k$-partite
graphon $T_k$,
\begin{equation}
 p=1-\frac1k,
 \qquad
 t(H,T_k)=\frac{\chi_H(k)}{k^{v(H)}}=\Phi_H(p).
 \label{eq:turan-equality}
\end{equation}
This equality motivates the following question.

\begin{definition}
\label{def:positive}
Let $H$ be non-bipartite.  We call $H$ \emph{positive} when every graphon $W$
of edge density
\begin{equation}
 p\in I_H:=\left[1-\frac1{\chi(H)-1},1\right]
 \label{eq:required-interval}
\end{equation}
satisfies $t(H,W)\geq\Phi_H(p)$.  We call $H$ \emph{negative} when such a
graphon and such a $p$ violate the inequality.
\end{definition}

The finite scope considered here consists of the NetworkX Graph Atlas
isomorphism classes on at most six vertices that are non-bipartite and have
neither isolated vertices nor an isolated-edge component.  The exclusions
remove factors that do not belong to the connected non-bipartite core of the
question.  A direct graph-atlas enumeration gives $117$ classes.

\begin{theorem}[Six-vertex classification]
\label{thm:classification}
Among the $117$ classes in the stated scope, exactly $94$ are positive and
$23$ are negative.  The negative Atlas indices are
\begin{equation}
\begin{split}
 48,50,129,140,141,143,149,152,158,159,166,170,172,173,{}\\
 182,184,186,189,190,197,198,204,206.
\end{split}
\label{eq:negative-list}
\end{equation}
Every other class in \Cref{tab:classification} is positive.
\end{theorem}

The theorem has two independently checkable parts.  For a negative row, a
finite rational sum evaluates both sides of the proposed inequality.  For a
positive row, the supplementary Lean theorem quantifies over arbitrary
probability spaces and graphons.  The proof is therefore not restricted to
finite graphs or step graphons.

The reusable mathematics is more informative than the row count.  Nineteen
positive rows follow from one pure-chordal theorem.  Five arise from a theorem
for cones over forests, and many others follow from rooted attachment or
closure arguments.  At the remaining rows, exact certificates replace a
large numerical search object by rational identities designed for formal
checking.  On the negative side, nineteen witnesses have one common tensor
construction; only four need local perturbations.

\paragraph{Proof archive.}
The supplementary archive contains the Lean development, exact certificate
files, rational checking scripts, and the generated tables included below.
Paths in \Cref{tab:row-sources} are relative to that archive.  The catalogue
table is regenerated by
\path{math_paper/scripts/generate_classification.py}; the negative
arithmetic is regenerated by
\path{math_paper/scripts/verify_negative_witnesses.py}.

\section{The chromatic target}

Write $n=v(H)$ and expand the chromatic polynomial as
$\chi_H(x)=\sum_{j=0}^n c_jx^j$.  Then
\begin{equation}
 \Phi_H(p)=\sum_{j=0}^n c_j(1-p)^{n-j},
 \label{eq:target-polynomial}
\end{equation}
so $\Phi_H$ is a polynomial in $p$ and its endpoint value is unambiguous.
For a clique, put
\begin{equation}
 A_r(p):=\Phi_{K_r}(p)
 =\prod_{i=1}^{r-1}\bigl(ip-(i-1)\bigr).
 \label{eq:clique-target}
\end{equation}
In particular, $A_2(p)=p$, $A_3(p)=p(2p-1)$, and
$A_4(p)=p(2p-1)(3p-2)$.

The target is compatible with several graph operations.  If $H=H_1\sqcup
H_2$, then
\begin{equation}
 t(H,W)=t(H_1,W)t(H_2,W),
 \qquad
 \Phi_H(p)=\Phi_{H_1}(p)\Phi_{H_2}(p).
 \label{eq:component-product}
\end{equation}
If $H$ is made from $N$ copies of a graph $F$ by identifying one chosen
vertex in every copy, then
\begin{equation}
 \chi_H(x)=\frac{\chi_F(x)^N}{x^{N-1}},
 \qquad
 \Phi_H(p)=\Phi_F(p)^N.
 \label{eq:vertex-amalgam-target}
\end{equation}
If instead the copies share one chosen edge, then
\begin{equation}
 \chi_H(x)=\frac{\chi_F(x)^N}{(x(x-1))^{N-1}},
 \qquad
 \Phi_H(p)=\frac{\Phi_F(p)^N}{p^{N-1}}.
 \label{eq:edge-amalgam-target}
\end{equation}
These identities follow by fixing the colours on the shared vertex or edge
and then colouring the copies independently.

\section{Analytic mechanisms for positive rows}

This section gives the common inequalities used in the positive
classification.  The specialized row arguments are indexed in
\Cref{tab:row-sources}; their Lean counterparts are identified there as
well.

\subsection{A weighted rooted triangle inequality}

Set
\begin{equation}
 d(x)=\int_\Omega W(x,y)\dd\mu(y),
 \qquad
 \tau(x)=\int_{\Omega^2}W(x,y)W(x,z)W(y,z)
 \dd\mu(y)\dd\mu(z).
 \label{eq:rooted-quantities}
\end{equation}

\begin{lemma}[Weighted rooted Goodman inequality]
\label{lem:weighted-goodman}
If $p>0$ and $h$ is a nonnegative integer, then
\begin{equation}
 \int_\Omega d(x)^h\tau(x)\dd\mu(x)
 \geq \frac{2p-1}{p}\int_\Omega d(x)^{h+2}\dd\mu(x).
 \label{eq:weighted-goodman}
\end{equation}
\end{lemma}

\begin{proof}
Put $M_j=\int d(x)^j\dd\mu(x)$, let $T_W$ be the integral operator with
kernel $W$, and define
\[
 I_h=\int d^h\tau,
 \qquad
 J_h=\ip{d^h}{T_Wd}.
\]
The inequality $ab\geq a+b-1$ on $[0,1]^2$, applied to $W(x,y)$ and
$W(x,z)$ and then multiplied by $d(x)^hW(y,z)$, gives
\begin{equation}
 I_h\geq2J_h-pM_h.
 \label{eq:I-J}
\end{equation}
Let $U=1-W$, $u=T_U\one=1-d$, and $q=1-p$.  Since
$T_Wd=d-q\one+T_Uu$,
\begin{equation}
 J_h=M_{h+1}-qM_h+\ip{d^h}{T_Uu}.
 \label{eq:J-decomposition}
\end{equation}
Symmetry of $U$ shows that twice
\[
 \ip{d^h}{T_Uu}-\int d(x)^hu(x)^2\dd\mu(x)
\]
equals
\[
 \int_{\Omega^2}U(x,y)(u(y)-u(x))(d(x)^h-d(y)^h)
 \dd\mu(x)\dd\mu(y).
\]
This is nonnegative because $d=1-u$.  Hence
\[
 J_h\geq M_{h+2}-M_{h+1}+pM_h.
\]
Using this in \eqref{eq:I-J} and subtracting the right-hand side of
\eqref{eq:weighted-goodman} gives
\[
 \frac1p\left(M_{h+2}-2pM_{h+1}+p^2M_h\right)
 =\frac1p\int d(x)^h(d(x)-p)^2\dd\mu(x)\geq0.
\]
\end{proof}

This lemma supplies the rooted triangle--tree, book, page, broom, and several
leaf-attachment rows in the catalogue.  A second local form is also useful.
With $B(x)=(T_Wd)(x)$,
\begin{equation}
 \tau(x)\geq\max\{2B(x)-p,0\}.
 \label{eq:pointwise-goodman}
\end{equation}
Indeed, the same inequality $ab\geq a+b-1$, now integrated only in the two
free variables, yields $\tau(x)\geq2B(x)-p$; nonnegativity supplies the other
branch.

\subsection{Cones over forests}

For $d(x)>0$, let $W_x$ be $W$ on the probability measure
$d\mu_x(y)=W(x,y)d\mu(y)/d(x)$, and let
$z_x=t(K_2,W_x)=\tau(x)/d(x)^2$.  At $d(x)=0$ set $z_x=0$.

\begin{theorem}[Cones over forests]
\label{thm:forest-cone}
Let $F$ be a forest with $n$ vertices, $e\geq1$ edges, and $c$ connected
components.  If $H=K_1\join F$ and $p\geq1/2$, then
\begin{equation}
 t(H,W)\geq p^c(2p-1)^e=\Phi_H(p).
 \label{eq:forest-cone}
\end{equation}
\end{theorem}

\begin{proof}
The forest case of Sidorenko's inequality
\cite[Section~5.2]{Lovasz2012} gives $t(F,W_x)\geq z_x^e$.
Conditioning on the image of the universal vertex gives
\begin{equation}
 t(H,W)=\int_\Omega d(x)^nt(F,W_x)\dd\mu(x).
 \label{eq:cone-link}
\end{equation}
By \Cref{lem:weighted-goodman},
\[
 \frac{\int d(x)^nz_x\dd\mu(x)}{\int d(x)^n\dd\mu(x)}
 \geq2-\frac1p.
\]
Apply Jensen's inequality first to $z\mapsto z^e$ with weight $d^n$ and
then to $d\mapsto d^n$.  Since a forest has $n=c+e$,
\[
 t(H,W)\geq p^n\left(2-\frac1p\right)^e
 =p^c(2p-1)^e.
\]
The chromatic-polynomial identity follows by colouring the cone vertex first
and then the forest.
\end{proof}

This theorem proves Atlas $40,111,117,135,$ and $136$.  Closely related
normalized-link arguments prove Atlas $49$, $156$, and $165$.  The six cone
lifts are Atlas $177$, $179$, $183$, $192$, $200$, and $205$.  Atlas $196$
follows by applying a
normalized-link cone theorem to the exact Atlas $43$ inequality.

\subsection{Odd cycles and pure chordal graphs}

\begin{theorem}[Odd-cycle bound]
\label{thm:odd-cycle}
For every odd $m\geq3$ and every graphon of edge density $p$,
\begin{equation}
 t(C_m,W)\geq p^m-p(1-p)^{m-1}=\Phi_{C_m}(p).
 \label{eq:odd-cycle}
\end{equation}
\end{theorem}

The proof separates $m=3,5,7$ from $m\geq9$.  For the three small values,
Goodman's triangle inequality~\cite{Goodman1959} is the case $m=3$.
For the three small values, inclusion--exclusion in $U=1-W$ reduces the inequality to sums of squares in
path moments.  For $m\geq9$, the integral operator of $U$ is split into its
constant direction and the mean-zero subspace.  The resulting spectral
moments are bounded separately on $p\geq2/3$ and $1/2<p<2/3$.  This is the
proof formalized for Atlas $7$ and $38$; it is independent of the earlier
odd-cycle project.  The full derivation is supplied with the archive and is
connected to these two row theorems in \Cref{tab:row-sources}.

The second large family is chordal.  For $r\geq3$, a \emph{pure
$r$-chordal graph} is a chordal graph whose maximal cliques all have order
$r$.

\begin{theorem}[Pure chordal bound]
\label{thm:pure-chordal}
If $H$ is pure $r$-chordal, then for every graphon $W$ with
$p\geq1-1/(r-1)$,
\begin{equation}
 t(H,W)\geq\Phi_H(p).
 \label{eq:pure-chordal}
\end{equation}
\end{theorem}

A clique-tree decomposition writes $H$ as $K_r$ bags joined along clique
separators.  Entropy gluing bounds the density of the resulting junction by
the bag densities divided by the separator densities.  The graphon
Moon and Moser's clique-ratio recurrence~\cite{MoonMoser1962} compares
consecutive clique densities, while the
chromatic polynomial of the same clique tree factors into the corresponding
$A_r/A_s$ terms from \eqref{eq:clique-target}; see also Agnarsson's
factorization for chordal graphs~\cite{Agnarsson2003}.  Multiplying these comparisons
proves \eqref{eq:pure-chordal}.  A regularization step removes divisions by
zero.  This proof gives the nineteen pure-chordal rows listed in
\Cref{tab:classification} and does not require connectedness.

\subsection{Closure and row-specific arguments}

Three closure rules reduce further rows to smaller ones.

\begin{proposition}[Closure rules]
\label{prop:closures}
Suppose all target functions below are nonnegative on the interval under
consideration.
\begin{enumerate}[label=(\roman*)]
\item Positivity is preserved under disjoint union.
\item If $H$ is a vertex self-amalgam of $N$ copies of $F$, then
$t(H,W)\geq t(F,W)^N$, and \eqref{eq:vertex-amalgam-target} transfers the
bound from $F$ to $H$.
\item If $H$ is an edge self-amalgam of $N$ copies of $F$, then a rooted
Cauchy--Schwarz argument gives
$t(H,W)\geq t(F,W)^N/p^{N-1}$, matching
\eqref{eq:edge-amalgam-target}.
\item Let $\operatorname{Wh}(F)$ attach one leaf to every vertex of $F$.
If $\Phi_F$ is nondecreasing on the required interval, then
\[
 t(\operatorname{Wh}(F),W)\geq p^{v(F)}\Phi_F(p)
 =\Phi_{\operatorname{Wh}(F)}(p).
\]
\end{enumerate}
\end{proposition}

The proof of each item conditions on the common rooted variables and applies
H\"older, Cauchy--Schwarz, or Jensen as appropriate.  In the whiskering
case, conditioning on the original vertices inserts one degree factor per
vertex; a change of measure followed by monotonicity of $\Phi_F$ yields the
stated bound.

The remaining analytic rows use a short list of two-dimensional reductions:
supporting planes for rooted density profiles, an exact square for Atlas
$145$, a Hilbert-space projection for Atlas $148$, and a concrete
clique--odd-cycle inequality for Atlas $187$.  The Atlas $126$ and $148$
arguments use the triangle-density profile proved by Fisher~\cite{Fisher1989}
and subsumed by Reiher's clique-density theorem~\cite{Reiher2016}.  The
triangle three-edge-tail argument uses the odd-walk inequality of Blekherman
and Raymond~\cite{BlekhermanRaymond2023}.  Each reduction ends in a
univariate polynomial inequality on the required interval.  The polynomial
steps are checked with rational factorization or Bernstein coefficients.
The row-to-proof index in \Cref{tab:row-sources} distinguishes these methods
from the exact rooted certificates in the next section.

\section{Exact rooted certificates}

The difficult positive rows are proved by rational identities over rooted
graph densities.  We briefly describe the common semantics.

Choose labelled root vertices and sample their images in $\Omega$.  Each
possible edge is represented by a Bernoulli variable whose conditional mean
is the corresponding value of $W$.  Repeated occurrences of the same
labelled edge use the same Bernoulli variable.  This convention is essential:
using independent copies would replace a factor $W$ by $W^2$ and would not
represent the graph density.

For a finite vector $v$ of rooted expressions and a rational positive
semidefinite matrix $Q$, the conditional expectation of $v^{\mathsf T}Qv$
is nonnegative.  An interval certificate combines such terms with the
nonnegative factors $p-a$ and $b-p$:
\begin{equation}
 t(H,W)-\Phi_H(p)
 =\mathbb E\!\left[v_0^{\mathsf T}Q_0v_0
 +(p-a)v_1^{\mathsf T}Q_1v_1
 +(b-p)v_2^{\mathsf T}Q_2v_2+R\right],
 \label{eq:interval-sos}
\end{equation}
where $R$ is a rational linear combination already known to be nonnegative.
Every equality in \eqref{eq:interval-sos} is checked coefficientwise, and
each $Q_i$ is supplied by a rational Gram factorization or a strict diagonal
dominance proof.

\begin{proposition}[Certificate rows]
\label{prop:certificate-rows}
The following exact certificates prove the indicated rows throughout their
required intervals.
\begin{enumerate}[label=(\roman*)]
\item Atlas $43$: a three-root identity on $[1/2,1]$ with $91$ nonzero
coefficients and two rational Gram corrections of orders $107$ and $48$.
\item Atlas
$118,122,124,147,151,153,157,168,169,171,174,181,185,188,194,199$:
four-root interval identities, each with $407$ coefficient equations.
\item Atlas $130$ on $[1/2,2/3]$ and Atlas $203$ on $[2/3,3/4]$:
induced rooted identities using $28$ positive integer Gram factors and
$936$ coefficient equations.  Four-root identities cover $[2/3,1]$ for
Atlas $130$ and $[3/4,1]$ for Atlas $203$.
\item Atlas $127$: a direct ground-$6$ split certificate on $[1/2,1]$.
\end{enumerate}
\end{proposition}

The direct ground-$6$ certificate in part (iv) is the proof used for Atlas
$127$ in this classification.  A stronger smoothed-Goodman statement was
also investigated, but it is not needed for the row theorem and is not used
here.  Atlas $43$ enters Atlas $196$ through the normalized-link implication
mentioned after \Cref{thm:forest-cone}.

The numerical searches that suggested these identities are not proof
inputs.  The final files contain rational coefficients, exact polynomial
equalities, and nonnegativity decompositions.  The Lean development connects
their rooted semantics to \eqref{eq:hom-density}.

\section{Exact counterexamples}

Let $\boldsymbol\mu=(\mu_1,\ldots,\mu_s)$ be rational nonnegative masses
summing to one, and let $A\in[0,1]^{s\times s}$ be a symmetric rational
matrix.  The associated step graphon satisfies
\begin{align}
 p&=\sum_{i,j=1}^s\mu_i\mu_jA_{ij},
 \label{eq:step-edge-density}\\
 t(H,W)&=\sum_{f:V(H)\to[s]}
 \prod_{u\in V(H)}\mu_{f(u)}
 \prod_{uv\in E(H)}A_{f(u),f(v)}.
 \label{eq:step-hom-density}
\end{align}
Thus every counterexample in \Cref{tab:negative-witnesses} is a finite exact
rational computation.

\subsection{Tensor products of Tur\'{a}n graphons}

On $[a]\times[b]$ with the uniform measure, define
\begin{equation}
 (T_a\otimes T_b)((i,j),(i',j'))
 =\one_{i\ne i'}\one_{j\ne j'}.
 \label{eq:tensor-turan}
\end{equation}

\begin{lemma}[Tensor evaluation]
\label{lem:tensor-evaluation}
For every finite graph $H$,
\begin{equation}
 p=\left(1-\frac1a\right)\left(1-\frac1b\right),
 \qquad
 t(H,T_a\otimes T_b)=
 \frac{\chi_H(a)\chi_H(b)}{(ab)^{v(H)}}.
 \label{eq:tensor-evaluation}
\end{equation}
\end{lemma}

\begin{proof}
An edge is present precisely when its two coordinates differ in both
positions, which gives the formula for $p$.  A homomorphism to the step graph
is exactly an ordered pair of proper colourings, one with $a$ colours and one
with $b$ colours.  There are $\chi_H(a)\chi_H(b)$ such pairs among
$(ab)^{v(H)}$ maps.
\end{proof}

Substitution into \eqref{eq:tensor-evaluation} proves strict failure for
nineteen of the rows in \eqref{eq:negative-list}.  The exact values appear in
\Cref{tab:negative-witnesses}.

\subsection{Four local witnesses}

Atlas $166$ and $172$ use uniform masses on four parts and
\begin{equation}
 A=\frac14
 \begin{pmatrix}
 0&1&4&4\\1&0&4&4\\4&4&0&1\\4&4&1&0
 \end{pmatrix}.
 \label{eq:two-scale-matrix}
\end{equation}
Atlas $206$ uses uniform masses on three parts and
\begin{equation}
 A=\frac14
 \begin{pmatrix}
 1&4&4\\4&0&4\\4&4&0
 \end{pmatrix}.
 \label{eq:one-diagonal-matrix}
\end{equation}
These give $p=9/16$ and $p=25/36$, respectively.

For Atlas $152$, use the five masses
\begin{equation}
 \boldsymbol\mu=\frac1{3750000}
 (1870752,2998,1875000,125,1125)
 \label{eq:atlas152-masses}
\end{equation}
and
\begin{equation}
 A=\frac1{16}
 \begin{pmatrix}
 0&0&16&16&0\\
 0&0&16&0&16\\
 16&16&0&1&16\\
 16&0&1&0&0\\
 0&16&16&0&0
 \end{pmatrix}.
 \label{eq:atlas152-matrix}
\end{equation}
Equations \eqref{eq:step-edge-density}--\eqref{eq:step-hom-density} give the
strict rational gaps in \Cref{tab:negative-witnesses}.  This finishes the
negative half of \Cref{thm:classification}.

\section{Formal verification}

The Lean definition of a graphon uses an arbitrary probability space and a
symmetric measurable kernel with values in $[0,1]$.  For each Atlas row, the
example module proves either \texttt{SatisfiesLowerBound} or
\texttt{ViolatesLowerBound}.  A source scan finds all $117$ modules marked
\texttt{.verified} and no occurrence of \texttt{sorry} in those modules.

The formal proof is intentionally row-oriented.  Some analytic notes state a
more general theorem, while Lean proves only the specialization needed for a
catalogue row.  The pure-chordal, rooted-SOS, and common negative constructions
have reusable method-level formalizations; several specialized analytic rows
instead expose only their final row theorem.  \Cref{tab:row-sources} names the
mathematical and Lean source for every row, so these scopes can be checked
without inferring a general theorem from a row specialization.

Twenty compact or specialized rows have retained sequential-build
measurements, reproduced in \Cref{tab:verification-cost}.  Their combined
time was $701.5$ minutes and the largest measured peak memory was
$11.91$ GiB.  Atlas $43$ used a separate $1097$-module sequential check:
$6.02$ hours and $15.4$ GiB peak memory.  A parallel attempt on a $64$ GiB
machine exhausted memory.  These measurements concern proof checking, not
the mathematical truth of the identities.

The author reports a successful full verification of the complete
development on another device.  The directory from that run and its exact
start and end times are unavailable in the present archive.  We therefore
separate that author-attested full build from the retained row-level
measurements and from the reproducible source scan.

\section{Proof of the classification theorem}

\begin{proof}[Proof of \Cref{thm:classification}]
The generation script enumerates the Graph Atlas, applies the stated scope
conditions, and independently checks each graph6 identifier, vertex count,
edge count, chromatic number, and chromatic polynomial.  It obtains $117$
rows.

For the $23$ indices in \eqref{eq:negative-list}, the rational checker uses
\Cref{lem:tensor-evaluation} or
\eqref{eq:step-edge-density}--\eqref{eq:atlas152-matrix} and verifies
$t(H,W)-\Phi_H(p)<0$.  All witness densities lie in $I_H$.

For each remaining row, \Cref{tab:classification} assigns one of the
positive methods in Sections 3--4, and \Cref{tab:row-sources} identifies the
corresponding mathematical source and Lean theorem.  The generated audit
checks that each of these $94$ modules has theorem kind
\texttt{SatisfiesLowerBound}; the formal theorem unfolds to the inequality in
\Cref{def:positive} for every graphon and every $p\in I_H$.  The two sets are
disjoint and contain all $117$ rows.  Hence $94$ rows are positive and $23$
are negative.
\end{proof}

\section{Discussion}

The classification has three mathematical features.  First, most positive
rows do not need case-specific semidefinite searches: clique-tree gluing,
rooted Goodman inequalities, links, and closure operations cover broad
families.  Second, the hard positive cases admit compact exact replacements
for much larger search outputs.  Third, the negative side is even more
uniform: tensor products of Tur\'{a}n graphons settle nineteen of twenty-three
rows.

The order-six boundary should not be interpreted as a claim about all finite
graphs.  Nor does the present work supply a simple graph-theoretic
characterization of positivity.  Its conclusion is a complete finite
classification together with reusable sufficient conditions and explicit
obstructions.  The machine-checked artifacts also make clear which claims are
general graphon theorems and which are formalized only at a fixed Atlas row.

\paragraph{AI assistance disclosure.}
OpenAI Codex drafted this manuscript from the project sources and generated
some of the organizational prose.  The mathematical claims are tied to the
exact scripts, certificates, and Lean sources described above.  The authors
remain responsible for the submitted text and conclusions.  A separate
companion article studies the use of ChatGPT Pro, Codex, and Claude Code in
the proof campaign~\cite{CompanionAI2026}.

\appendix

\section{Complete classification}
\label{app:classification}
\begin{landscape}
\fontsize{5.5}{6.5}\selectfont
\setlength{\tabcolsep}{1.5pt}
\begin{longtable}{r l r r r l p{4.0cm} p{3.5cm} p{2.7cm}}
\caption{Complete classification.  The interval column gives the required range of edge densities.}\label{tab:classification}\\
\toprule
Atlas & graph6 & $v$ & $e$ & $\chi$ & interval & $\Phi_H(p)$ & method & Lean theorem\\
\midrule
\endfirsthead
\multicolumn{9}{c}{\tablename\ \thetable\ (continued)}\\
\toprule
Atlas & graph6 & $v$ & $e$ & $\chi$ & interval & $\Phi_H(p)$ & method & Lean theorem\\
\midrule
\endhead
\midrule
\multicolumn{9}{r}{continued on next page}\\
\endfoot
\bottomrule
\endlastfoot
7 & \texttt{Bw} & 3 & 3 & 3 & $[\frac{1}{2},1]$ & $p \left(2 p - 1\right)$ & + odd-cycle bound & \texttt{Graph007.status}\\
15 & \texttt{CN} & 4 & 4 & 3 & $[\frac{1}{2},1]$ & $p^{2} \cdot \left(2 p - 1\right)$ & + paw inequality & \texttt{Graph015.status}\\
17 & \texttt{C|} & 4 & 5 & 3 & $[\frac{1}{2},1]$ & $p \left(2 p - 1\right)^{2}$ & + pure chordal & \texttt{Graph017.status}\\
18 & \texttt{C\textasciitilde{}} & 4 & 6 & 4 & $[\frac{2}{3},1]$ & $p \left(2 p - 1\right) \left(3 p - 2\right)$ & + pure chordal & \texttt{Graph018.status}\\
34 & \texttt{D@\{} & 5 & 5 & 3 & $[\frac{1}{2},1]$ & $p^{3} \cdot \left(2 p - 1\right)$ & + rooted triangle--tree inequality & \texttt{Graph034.status}\\
35 & \texttt{Dx\_} & 5 & 5 & 3 & $[\frac{1}{2},1]$ & $p^{3} \cdot \left(2 p - 1\right)$ & + two-root book inequality & \texttt{Graph035.status}\\
36 & \texttt{DJc} & 5 & 5 & 3 & $[\frac{1}{2},1]$ & $p^{3} \cdot \left(2 p - 1\right)$ & + rooted triangle--tree inequality & \texttt{Graph036.status}\\
38 & \texttt{Dhc} & 5 & 5 & 3 & $[\frac{1}{2},1]$ & $p \left(2 p - 1\right) \left(2 p^{2} - 2 p + 1\right)$ & + odd-cycle bound & \texttt{Graph038.status}\\
40 & \texttt{DjW} & 5 & 6 & 3 & $[\frac{1}{2},1]$ & $p^{2} \left(2 p - 1\right)^{2}$ & + cone over a forest & \texttt{Graph040.status}\\
41 & \texttt{Db[} & 5 & 6 & 3 & $[\frac{1}{2},1]$ & $p^{2} \left(2 p - 1\right)^{2}$ & + rooted triangle--tree inequality & \texttt{Graph041.status}\\
42 & \texttt{D`\{} & 5 & 6 & 3 & $[\frac{1}{2},1]$ & $p^{2} \left(2 p - 1\right)^{2}$ & + pure chordal & \texttt{Graph042.status}\\
43 & \texttt{Dlc} & 5 & 6 & 3 & $[\frac{1}{2},1]$ & $p \left(2 p - 1\right) \left(3 p^{2} - 3 p + 1\right)$ & + three-root interval SOS & \texttt{Graph043.status}\\
45 & \texttt{DJ\{} & 5 & 7 & 4 & $[\frac{2}{3},1]$ & $p^{2} \cdot \left(2 p - 1\right) \left(3 p - 2\right)$ & + clique common-leaf inequality & \texttt{Graph045.status}\\
46 & \texttt{DF\{} & 5 & 7 & 3 & $[\frac{1}{2},1]$ & $p \left(2 p - 1\right)^{3}$ & + pure chordal & \texttt{Graph046.status}\\
47 & \texttt{Djs} & 5 & 7 & 3 & $[\frac{1}{2},1]$ & $p \left(2 p - 1\right)^{3}$ & + pure chordal & \texttt{Graph047.status}\\
48 & \texttt{D]w} & 5 & 7 & 3 & $[\frac{1}{2},1]$ & $p \left(2 p - 1\right) \left(5 p^{2} - 6 p + 2\right)$ & -- tensor product of Turan graphons & \texttt{Graph048.status}\\
49 & \texttt{Df\{} & 5 & 8 & 4 & $[\frac{2}{3},1]$ & $p \left(2 p - 1\right)^{2} \cdot \left(3 p - 2\right)$ & + special cone inequality & \texttt{Graph049.status}\\
50 & \texttt{Dl\{} & 5 & 8 & 3 & $[\frac{1}{2},1]$ & $p \left(2 p - 1\right) \left(7 p^{2} - 9 p + 3\right)$ & -- tensor product of Turan graphons & \texttt{Graph050.status}\\
51 & \texttt{Dn\{} & 5 & 9 & 4 & $[\frac{2}{3},1]$ & $p \left(2 p - 1\right) \left(3 p - 2\right)^{2}$ & + pure chordal & \texttt{Graph051.status}\\
52 & \texttt{D\textasciitilde{}\{} & 5 & 10 & 5 & $[\frac{3}{4},1]$ & $p \left(2 p - 1\right) \left(3 p - 2\right) \left(4 p - 3\right)$ & + pure chordal & \texttt{Graph052.status}\\
84 & \texttt{ECd\_} & 6 & 5 & 3 & $[\frac{1}{2},1]$ & $p^{3} \cdot \left(2 p - 1\right)$ & + component multiplicativity & \texttt{Graph084.status}\\
92 & \texttt{E?Fw} & 6 & 6 & 3 & $[\frac{1}{2},1]$ & $p^{4} \cdot \left(2 p - 1\right)$ & + rooted triangle--tree inequality & \texttt{Graph092.status}\\
93 & \texttt{EC\{O} & 6 & 6 & 3 & $[\frac{1}{2},1]$ & $p^{4} \cdot \left(2 p - 1\right)$ & + two-root book inequality & \texttt{Graph093.status}\\
94 & \texttt{E@dW} & 6 & 6 & 3 & $[\frac{1}{2},1]$ & $p^{4} \cdot \left(2 p - 1\right)$ & + whiskering & \texttt{Graph094.status}\\
95 & \texttt{EG\}?} & 6 & 6 & 3 & $[\frac{1}{2},1]$ & $p^{4} \cdot \left(2 p - 1\right)$ & + mixed rooted-triangle inequality & \texttt{Graph095.status}\\
97 & \texttt{EYWO} & 6 & 6 & 3 & $[\frac{1}{2},1]$ & $p^{4} \cdot \left(2 p - 1\right)$ & + adjacent leaf--tail inequality & \texttt{Graph097.status}\\
100 & \texttt{EsCW} & 6 & 6 & 3 & $[\frac{1}{2},1]$ & $p^{4} \cdot \left(2 p - 1\right)$ & + triangle broom inequality & \texttt{Graph100.status}\\
102 & \texttt{EJe?} & 6 & 6 & 3 & $[\frac{1}{2},1]$ & $p^{4} \cdot \left(2 p - 1\right)$ & + triangle tail inequality & \texttt{Graph102.status}\\
104 & \texttt{Ehd?} & 6 & 6 & 3 & $[\frac{1}{2},1]$ & $p^{2} \cdot \left(2 p - 1\right) \left(2 p^{2} - 2 p + 1\right)$ & + odd-cycle leaf inequality & \texttt{Graph104.status}\\
106 & \texttt{EJaG} & 6 & 6 & 3 & $[\frac{1}{2},1]$ & $p^{2} \left(2 p - 1\right)^{2}$ & + component multiplicativity & \texttt{Graph106.status}\\
111 & \texttt{EB\{G} & 6 & 7 & 3 & $[\frac{1}{2},1]$ & $p^{3} \left(2 p - 1\right)^{2}$ & + cone over a forest & \texttt{Graph111.status}\\
112 & \texttt{EhX\_} & 6 & 7 & 3 & $[\frac{1}{2},1]$ & $p^{3} \left(2 p - 1\right)^{2}$ & + two-root book inequality & \texttt{Graph112.status}\\
113 & \texttt{E\textasciicircum{}\_O} & 6 & 7 & 3 & $[\frac{1}{2},1]$ & $p^{3} \left(2 p - 1\right)^{2}$ & + diamond leaf inequality & \texttt{Graph113.status}\\
114 & \texttt{EJwG} & 6 & 7 & 3 & $[\frac{1}{2},1]$ & $p^{3} \left(2 p - 1\right)^{2}$ & + rooted triangle--tree inequality & \texttt{Graph114.status}\\
115 & \texttt{E`Xg} & 6 & 7 & 3 & $[\frac{1}{2},1]$ & $p^{3} \left(2 p - 1\right)^{2}$ & + self-amalgamation & \texttt{Graph115.status}\\
117 & \texttt{EtaG} & 6 & 7 & 3 & $[\frac{1}{2},1]$ & $p^{3} \left(2 p - 1\right)^{2}$ & + cone over a forest & \texttt{Graph117.status}\\
118 & \texttt{Eht?} & 6 & 7 & 3 & $[\frac{1}{2},1]$ & $p^{2} \cdot \left(2 p - 1\right) \left(3 p^{2} - 3 p + 1\right)$ & + Compact four-root SOS & \texttt{Graph118.status}\\
119 & \texttt{EtoO} & 6 & 7 & 3 & $[\frac{1}{2},1]$ & $p^{3} \left(2 p - 1\right)^{2}$ & + bowtie leaf inequality & \texttt{Graph119.status}\\
120 & \texttt{EB\}?} & 6 & 7 & 3 & $[\frac{1}{2},1]$ & $p^{3} \left(2 p - 1\right)^{2}$ & + triangle-book attachment inequality & \texttt{Graph120.status}\\
122 & \texttt{Eld?} & 6 & 7 & 3 & $[\frac{1}{2},1]$ & $p^{2} \cdot \left(2 p - 1\right) \left(3 p^{2} - 3 p + 1\right)$ & + Compact four-root SOS & \texttt{Graph122.status}\\
123 & \texttt{EJy?} & 6 & 7 & 3 & $[\frac{1}{2},1]$ & $p^{3} \left(2 p - 1\right)^{2}$ & + triangle-book attachment inequality & \texttt{Graph123.status}\\
124 & \texttt{Exd?} & 6 & 7 & 3 & $[\frac{1}{2},1]$ & $p^{2} \cdot \left(2 p - 1\right) \left(3 p^{2} - 3 p + 1\right)$ & + Compact four-root SOS & \texttt{Graph124.status}\\
126 & \texttt{ERUO} & 6 & 7 & 3 & $[\frac{1}{2},1]$ & $p^{2} \cdot \left(2 p - 1\right) \left(3 p^{2} - 3 p + 1\right)$ & + Atlas 126 triangle--\$C\_4\$ vertex supporting-plane theorem & \texttt{Graph126.status}\\
127 & \texttt{EZEG} & 6 & 7 & 3 & $[\frac{1}{2},1]$ & $p \left(2 p - 1\right)^{2} \cdot \left(2 p^{2} - 2 p + 1\right)$ & + direct compact interval SOS & \texttt{Graph127.status}\\
129 & \texttt{EheO} & 6 & 7 & 3 & $[\frac{1}{2},1]$ & $p \left(2 p - 1\right) \left(5 p^{3} - 8 p^{2} + 5 p - 1\right)$ & -- tensor product of Turan graphons & \texttt{Graph129.status}\\
130 & \texttt{E\{CW} & 6 & 7 & 3 & $[\frac{1}{2},1]$ & $p^{3} \left(2 p - 1\right)^{2}$ & + Induced and four-root SOS & \texttt{Graph130.status}\\
133 & \texttt{E\textasciitilde{}a?} & 6 & 8 & 4 & $[\frac{2}{3},1]$ & $p^{3} \cdot \left(2 p - 1\right) \left(3 p - 2\right)$ & + clique common-leaf inequality & \texttt{Graph133.status}\\
134 & \texttt{E\textasciitilde{}\_O} & 6 & 8 & 4 & $[\frac{2}{3},1]$ & $p^{3} \cdot \left(2 p - 1\right) \left(3 p - 2\right)$ & + clique distributed-leaf inequality & \texttt{Graph134.status}\\
135 & \texttt{EzW\_} & 6 & 8 & 3 & $[\frac{1}{2},1]$ & $p^{2} \left(2 p - 1\right)^{3}$ & + cone over a forest & \texttt{Graph135.status}\\
136 & \texttt{Ejt?} & 6 & 8 & 3 & $[\frac{1}{2},1]$ & $p^{2} \left(2 p - 1\right)^{3}$ & + cone over a forest & \texttt{Graph136.status}\\
137 & \texttt{EjsG} & 6 & 8 & 3 & $[\frac{1}{2},1]$ & $p^{2} \left(2 p - 1\right)^{3}$ & + triangle-book attachment inequality & \texttt{Graph137.status}\\
138 & \texttt{Ez`\_} & 6 & 8 & 3 & $[\frac{1}{2},1]$ & $p^{2} \left(2 p - 1\right)^{3}$ & + rooted triangle--tree inequality & \texttt{Graph138.status}\\
139 & \texttt{Eju?} & 6 & 8 & 3 & $[\frac{1}{2},1]$ & $p^{2} \left(2 p - 1\right)^{3}$ & + triangle-book attachment inequality & \texttt{Graph139.status}\\
140 & \texttt{Ev`\_} & 6 & 8 & 3 & $[\frac{1}{2},1]$ & $p^{2} \cdot \left(2 p - 1\right) \left(5 p^{2} - 6 p + 2\right)$ & -- tensor product of Turan graphons & \texttt{Graph140.status}\\
141 & \texttt{EXSw} & 6 & 8 & 3 & $[\frac{1}{2},1]$ & $p^{2} \cdot \left(2 p - 1\right) \left(5 p^{2} - 6 p + 2\right)$ & -- tensor product of Turan graphons & \texttt{Graph141.status}\\
142 & \texttt{E\textasciitilde{}AG} & 6 & 8 & 4 & $[\frac{2}{3},1]$ & $p^{3} \cdot \left(2 p - 1\right) \left(3 p - 2\right)$ & + supporting-plane inequality & \texttt{Graph142.status}\\
143 & \texttt{Er`o} & 6 & 8 & 3 & $[\frac{1}{2},1]$ & $p^{2} \cdot \left(2 p - 1\right) \left(5 p^{2} - 6 p + 2\right)$ & -- tensor product of Turan graphons & \texttt{Graph143.status}\\
144 & \texttt{EB\}G} & 6 & 8 & 3 & $[\frac{1}{2},1]$ & $p^{2} \left(2 p - 1\right)^{3}$ & + pure chordal & \texttt{Graph144.status}\\
145 & \texttt{Exe\_} & 6 & 8 & 3 & $[\frac{1}{2},1]$ & $p \left(2 p - 1\right)^{2} \cdot \left(3 p^{2} - 3 p + 1\right)$ & + exact square identity & \texttt{Graph145.status}\\
147 & \texttt{EhMg} & 6 & 8 & 3 & $[\frac{1}{2},1]$ & $p \left(2 p - 1\right)^{2} \cdot \left(3 p^{2} - 3 p + 1\right)$ & + Compact four-root SOS & \texttt{Graph147.status}\\
148 & \texttt{EyUG} & 6 & 8 & 3 & $[\frac{1}{2},1]$ & $p \left(2 p - 1\right)^{2} \cdot \left(3 p^{2} - 3 p + 1\right)$ & + Hilbert projection & \texttt{Graph148.status}\\
149 & \texttt{Ele\_} & 6 & 8 & 3 & $[\frac{1}{2},1]$ & $p \left(2 p - 1\right) \left(7 p^{3} - 11 p^{2} + 6 p - 1\right)$ & -- tensor product of Turan graphons & \texttt{Graph149.status}\\
150 & \texttt{EJyG} & 6 & 8 & 3 & $[\frac{1}{2},1]$ & $p^{2} \left(2 p - 1\right)^{3}$ & + pure chordal & \texttt{Graph150.status}\\
151 & \texttt{EhdW} & 6 & 8 & 3 & $[\frac{1}{2},1]$ & $p \left(2 p - 1\right)^{2} \cdot \left(4 p^{2} - 5 p + 2\right)$ & + Compact four-root SOS & \texttt{Graph151.status}\\
152 & \texttt{EhNG} & 6 & 8 & 3 & $[\frac{1}{2},1]$ & $p \left(2 p - 1\right)^{2} \cdot \left(3 p^{2} - 3 p + 1\right)$ & -- fractional-fibre step graphon & \texttt{Graph152.status}\\
153 & \texttt{Ehf\_} & 6 & 8 & 3 & $[\frac{1}{2},1]$ & $p \left(2 p - 1\right) \left(7 p^{3} - 12 p^{2} + 8 p - 2\right)$ & + Compact four-root SOS & \texttt{Graph153.status}\\
156 & \texttt{E\textasciitilde{}H\_} & 6 & 9 & 4 & $[\frac{2}{3},1]$ & $p^{2} \left(2 p - 1\right)^{2} \cdot \left(3 p - 2\right)$ & + special cone inequality & \texttt{Graph156.status}\\
157 & \texttt{E\textasciitilde{}`\_} & 6 & 9 & 4 & $[\frac{2}{3},1]$ & $p^{2} \left(2 p - 1\right)^{2} \cdot \left(3 p - 2\right)$ & + Compact four-root SOS & \texttt{Graph157.status}\\
158 & \texttt{El\{G} & 6 & 9 & 3 & $[\frac{1}{2},1]$ & $p^{2} \cdot \left(2 p - 1\right) \left(7 p^{2} - 9 p + 3\right)$ & -- tensor product of Turan graphons & \texttt{Graph158.status}\\
159 & \texttt{EZSw} & 6 & 9 & 3 & $[\frac{1}{2},1]$ & $p^{2} \cdot \left(2 p - 1\right) \left(7 p^{2} - 9 p + 3\right)$ & -- tensor product of Turan graphons & \texttt{Graph159.status}\\
160 & \texttt{E\textasciitilde{}@g} & 6 & 9 & 4 & $[\frac{2}{3},1]$ & $p^{2} \left(2 p - 1\right)^{2} \cdot \left(3 p - 2\right)$ & + supporting-plane inequality & \texttt{Graph160.status}\\
161 & \texttt{E?\textasciitilde{}w} & 6 & 9 & 3 & $[\frac{1}{2},1]$ & $p \left(2 p - 1\right)^{4}$ & + pure chordal & \texttt{Graph161.status}\\
162 & \texttt{E|e\_} & 6 & 9 & 3 & $[\frac{1}{2},1]$ & $p \left(2 p - 1\right)^{4}$ & + pure chordal & \texttt{Graph162.status}\\
163 & \texttt{EyuG} & 6 & 9 & 3 & $[\frac{1}{2},1]$ & $p \left(2 p - 1\right)^{4}$ & + pure chordal & \texttt{Graph163.status}\\
164 & \texttt{EyVG} & 6 & 9 & 3 & $[\frac{1}{2},1]$ & $p \left(2 p - 1\right)^{4}$ & + pure chordal & \texttt{Graph164.status}\\
165 & \texttt{E\textasciitilde{}aG} & 6 & 9 & 4 & $[\frac{2}{3},1]$ & $p^{2} \left(2 p - 1\right)^{2} \cdot \left(3 p - 2\right)$ & + special cone inequality & \texttt{Graph165.status}\\
166 & \texttt{ElfO} & 6 & 9 & 3 & $[\frac{1}{2},1]$ & $p \left(2 p - 1\right)^{2} \cdot \left(5 p^{2} - 6 p + 2\right)$ & -- local perturbation of a Turan graphon & \texttt{Graph166.status}\\
167 & \texttt{E\textasciicircum{}eG} & 6 & 9 & 3 & $[\frac{1}{2},1]$ & $p \left(2 p - 1\right)^{4}$ & + pure chordal & \texttt{Graph167.status}\\
168 & \texttt{E\textasciicircum{}MG} & 6 & 9 & 3 & $[\frac{1}{2},1]$ & $p \left(2 p - 1\right)^{2} \cdot \left(5 p^{2} - 6 p + 2\right)$ & + Compact four-root SOS & \texttt{Graph168.status}\\
169 & \texttt{Exf\_} & 6 & 9 & 4 & $[\frac{2}{3},1]$ & $p \left(2 p - 1\right) \left(3 p - 2\right) \left(3 p^{2} - 3 p + 1\right)$ & + Compact four-root SOS & \texttt{Graph169.status}\\
170 & \texttt{EO\textasciitilde{}o} & 6 & 9 & 3 & $[\frac{1}{2},1]$ & $p \left(2 p - 1\right) \left(11 p^{3} - 19 p^{2} + 11 p - 2\right)$ & -- tensor product of Turan graphons & \texttt{Graph170.status}\\
171 & \texttt{Ehew} & 6 & 9 & 3 & $[\frac{1}{2},1]$ & $p \left(2 p - 1\right) \left(11 p^{3} - 20 p^{2} + 13 p - 3\right)$ & + Compact four-root SOS & \texttt{Graph171.status}\\
172 & \texttt{Elf\_} & 6 & 9 & 3 & $[\frac{1}{2},1]$ & $p \left(2 p - 1\right)^{2} \cdot \left(5 p^{2} - 6 p + 2\right)$ & -- local perturbation of a Turan graphon & \texttt{Graph172.status}\\
173 & \texttt{ElMg} & 6 & 9 & 3 & $[\frac{1}{2},1]$ & $p \left(2 p - 1\right)^{2} \cdot \left(6 p^{2} - 8 p + 3\right)$ & -- tensor product of Turan graphons & \texttt{Graph173.status}\\
174 & \texttt{EtTg} & 6 & 9 & 3 & $[\frac{1}{2},1]$ & $p \left(2 p - 1\right) \left(13 p^{3} - 25 p^{2} + 17 p - 4\right)$ & + Compact four-root SOS & \texttt{Graph174.status}\\
177 & \texttt{En\{G} & 6 & 10 & 4 & $[\frac{2}{3},1]$ & $p^{2} \cdot \left(2 p - 1\right) \left(3 p - 2\right)^{2}$ & + six explicit cone lifts & \texttt{Graph177.status}\\
178 & \texttt{En\}?} & 6 & 10 & 4 & $[\frac{2}{3},1]$ & $p^{2} \cdot \left(2 p - 1\right) \left(3 p - 2\right)^{2}$ & + supporting-plane inequality & \texttt{Graph178.status}\\
179 & \texttt{E\_\textasciitilde{}w} & 6 & 10 & 4 & $[\frac{2}{3},1]$ & $p \left(2 p - 1\right)^{3} \cdot \left(3 p - 2\right)$ & + six explicit cone lifts & \texttt{Graph179.status}\\
180 & \texttt{EjtW} & 6 & 10 & 4 & $[\frac{2}{3},1]$ & $p \left(2 p - 1\right)^{3} \cdot \left(3 p - 2\right)$ & + two-root book inequality & \texttt{Graph180.status}\\
181 & \texttt{E\textasciicircum{}mG} & 6 & 10 & 4 & $[\frac{2}{3},1]$ & $p \left(2 p - 1\right)^{3} \cdot \left(3 p - 2\right)$ & + Compact four-root SOS & \texttt{Graph181.status}\\
182 & \texttt{E\textasciicircum{}Mg} & 6 & 10 & 3 & $[\frac{1}{2},1]$ & $p \left(2 p - 1\right)^{2} \cdot \left(7 p^{2} - 9 p + 3\right)$ & -- tensor product of Turan graphons & \texttt{Graph182.status}\\
183 & \texttt{EjvG} & 6 & 10 & 4 & $[\frac{2}{3},1]$ & $p \left(2 p - 1\right)^{3} \cdot \left(3 p - 2\right)$ & + six explicit cone lifts & \texttt{Graph183.status}\\
184 & \texttt{Elfo} & 6 & 10 & 3 & $[\frac{1}{2},1]$ & $p \left(2 p - 1\right)^{2} \cdot \left(7 p^{2} - 9 p + 3\right)$ & -- tensor product of Turan graphons & \texttt{Graph184.status}\\
185 & \texttt{Exv\_} & 6 & 10 & 4 & $[\frac{2}{3},1]$ & $p \left(2 p - 1\right) \left(3 p - 2\right) \left(5 p^{2} - 6 p + 2\right)$ & + Compact four-root SOS & \texttt{Graph185.status}\\
186 & \texttt{ErXw} & 6 & 10 & 3 & $[\frac{1}{2},1]$ & $p \left(2 p - 1\right)^{2} \cdot \left(8 p^{2} - 11 p + 4\right)$ & -- tensor product of Turan graphons & \texttt{Graph186.status}\\
187 & \texttt{Ehfw} & 6 & 10 & 4 & $[\frac{2}{3},1]$ & $p \left(2 p - 1\right) \left(3 p - 2\right) \left(5 p^{2} - 6 p + 2\right)$ & + concrete clique--odd-cycle theorem & \texttt{Graph187.status}\\
188 & \texttt{EzNG} & 6 & 10 & 3 & $[\frac{1}{2},1]$ & $p \left(2 p - 1\right) \left(17 p^{3} - 33 p^{2} + 22 p - 5\right)$ & + Compact four-root SOS & \texttt{Graph188.status}\\
189 & \texttt{E\textasciicircum{}NG} & 6 & 10 & 3 & $[\frac{1}{2},1]$ & $p \left(2 p - 1\right) \left(19 p^{3} - 38 p^{2} + 26 p - 6\right)$ & -- tensor product of Turan graphons & \texttt{Graph189.status}\\
190 & \texttt{EyUw} & 6 & 10 & 3 & $[\frac{1}{2},1]$ & $p \left(2 p - 1\right)^{2} \cdot \left(8 p^{2} - 11 p + 4\right)$ & -- tensor product of Turan graphons & \texttt{Graph190.status}\\
191 & \texttt{E\textasciitilde{}|?} & 6 & 11 & 5 & $[\frac{3}{4},1]$ & $p^{2} \cdot \left(2 p - 1\right) \left(3 p - 2\right) \left(4 p - 3\right)$ & + clique common-leaf inequality & \texttt{Graph191.status}\\
192 & \texttt{E\textasciitilde{}Xo} & 6 & 11 & 4 & $[\frac{2}{3},1]$ & $p \left(2 p - 1\right)^{2} \left(3 p - 2\right)^{2}$ & + six explicit cone lifts & \texttt{Graph192.status}\\
193 & \texttt{En\}G} & 6 & 11 & 4 & $[\frac{2}{3},1]$ & $p \left(2 p - 1\right)^{2} \left(3 p - 2\right)^{2}$ & + rooted triangle--tree inequality & \texttt{Graph193.status}\\
194 & \texttt{E\textasciitilde{}wW} & 6 & 11 & 4 & $[\frac{2}{3},1]$ & $p \left(2 p - 1\right) \left(3 p - 2\right) \left(7 p^{2} - 9 p + 3\right)$ & + Compact four-root SOS & \texttt{Graph194.status}\\
195 & \texttt{EyVw} & 6 & 11 & 4 & $[\frac{2}{3},1]$ & $p \left(2 p - 1\right)^{2} \left(3 p - 2\right)^{2}$ & + pure chordal & \texttt{Graph195.status}\\
196 & \texttt{ER\textasciitilde{}g} & 6 & 11 & 4 & $[\frac{2}{3},1]$ & $p \left(2 p - 1\right) \left(3 p - 2\right) \left(7 p^{2} - 9 p + 3\right)$ & + normalized-link cone lift & \texttt{Graph196.status}\\
197 & \texttt{E\}\textasciicircum{}G} & 6 & 11 & 3 & $[\frac{1}{2},1]$ & $p \left(2 p - 1\right) \left(23 p^{3} - 46 p^{2} + 31 p - 7\right)$ & -- tensor product of Turan graphons & \texttt{Graph197.status}\\
198 & \texttt{Ep\textasciitilde{}o} & 6 & 11 & 3 & $[\frac{1}{2},1]$ & $p \left(2 p - 1\right) \left(23 p^{3} - 46 p^{2} + 31 p - 7\right)$ & -- tensor product of Turan graphons & \texttt{Graph198.status}\\
199 & \texttt{El\textasciicircum{}g} & 6 & 11 & 4 & $[\frac{2}{3},1]$ & $p \left(2 p - 1\right) \left(3 p - 2\right) \left(8 p^{2} - 11 p + 4\right)$ & + Compact four-root SOS & \texttt{Graph199.status}\\
200 & \texttt{E\textasciitilde{}\{W} & 6 & 12 & 5 & $[\frac{3}{4},1]$ & $p \left(2 p - 1\right)^{2} \cdot \left(3 p - 2\right) \left(4 p - 3\right)$ & + six explicit cone lifts & \texttt{Graph200.status}\\
201 & \texttt{E\textasciitilde{}z\_} & 6 & 12 & 4 & $[\frac{2}{3},1]$ & $p \left(2 p - 1\right) \left(3 p - 2\right)^{3}$ & + pure chordal & \texttt{Graph201.status}\\
202 & \texttt{Ep\textasciitilde{}w} & 6 & 12 & 4 & $[\frac{2}{3},1]$ & $p \left(2 p - 1\right) \left(3 p - 2\right)^{3}$ & + pure chordal & \texttt{Graph202.status}\\
203 & \texttt{E\textasciitilde{}\textasciicircum{}G} & 6 & 12 & 4 & $[\frac{2}{3},1]$ & $p \left(2 p - 1\right) \left(3 p - 2\right) \left(10 p^{2} - 14 p + 5\right)$ & + Induced and four-root SOS & \texttt{Graph203.status}\\
204 & \texttt{EznW} & 6 & 12 & 3 & $[\frac{1}{2},1]$ & $p \left(2 p - 1\right) \left(32 p^{3} - 67 p^{2} + 47 p - 11\right)$ & -- tensor product of Turan graphons & \texttt{Graph204.status}\\
205 & \texttt{E\textasciitilde{}\textasciitilde{}G} & 6 & 13 & 5 & $[\frac{3}{4},1]$ & $p \left(2 p - 1\right) \left(3 p - 2\right)^{2} \cdot \left(4 p - 3\right)$ & + six explicit cone lifts & \texttt{Graph205.status}\\
206 & \texttt{E\textasciitilde{}nW} & 6 & 13 & 4 & $[\frac{2}{3},1]$ & $p \left(2 p - 1\right) \left(3 p - 2\right) \left(13 p^{2} - 19 p + 7\right)$ & -- local perturbation of a Turan graphon & \texttt{Graph206.status}\\
207 & \texttt{Ez\textasciitilde{}w} & 6 & 14 & 5 & $[\frac{3}{4},1]$ & $p \left(2 p - 1\right) \left(3 p - 2\right) \left(4 p - 3\right)^{2}$ & + pure chordal & \texttt{Graph207.status}\\
208 & \texttt{E\textasciitilde{}\textasciitilde{}w} & 6 & 15 & 6 & $[\frac{4}{5},1]$ & $p \left(2 p - 1\right) \left(3 p - 2\right) \left(4 p - 3\right) \left(5 p - 4\right)$ & + pure chordal & \texttt{Graph208.status}\\
\end{longtable}
\end{landscape}


\section{Row-to-source index}
\label{app:row-sources}
\begin{landscape}
\scriptsize
\begin{longtable}{r l p{8.7cm} p{7.0cm}}
\caption{Machine-readable identifiers, edge lists, and source files.}\label{tab:row-sources}\\
\toprule Atlas & graph6 & edge list & source\\ \midrule
\endfirsthead
\toprule Atlas & graph6 & edge list & source\\ \midrule
\endhead
\bottomrule\endfoot
7 & \texttt{Bw} & 01, 02, 12 & \path{notes/paper_new_region2_v3.tex}; \path{lean/Taeyoung/Examples/Graph007.lean}\\
15 & \texttt{CN} & 03, 12, 13, 23 & \path{notes/rooted_triangle_tree_extensions.tex}; \path{lean/Taeyoung/Examples/Graph015.lean}\\
17 & \texttt{C|} & 01, 02, 03, 12, 23 & \path{notes/pure_chordal.tex}; \path{lean/Taeyoung/Examples/Graph017.lean}\\
18 & \texttt{C\textasciitilde{}} & 01, 02, 03, 12, 13, 23 & \path{notes/pure_chordal.tex}; \path{lean/Taeyoung/Examples/Graph018.lean}\\
34 & \texttt{D@\{} & 04, 14, 23, 24, 34 & \path{notes/rooted_triangle_tree_extensions.tex}; \path{lean/Taeyoung/Examples/Graph034.lean}\\
35 & \texttt{Dx\_} & 01, 02, 04, 12, 23 & \path{notes/two_root_triangle_leaf_cones.tex}; \path{lean/Taeyoung/Examples/Graph035.lean}\\
36 & \texttt{DJc} & 04, 12, 13, 23, 34 & \path{notes/rooted_triangle_tree_extensions.tex}; \path{lean/Taeyoung/Examples/Graph036.lean}\\
38 & \texttt{Dhc} & 01, 04, 12, 23, 34 & \path{notes/paper_new_region2_v3.tex}; \path{lean/Taeyoung/Examples/Graph038.lean}\\
40 & \texttt{DjW} & 01, 12, 13, 14, 23, 24 & \path{notes/forest_cone_graphon_bound.tex}; \path{lean/Taeyoung/Examples/Graph040.lean}\\
41 & \texttt{Db[} & 01, 13, 14, 23, 24, 34 & \path{notes/page_rooted_triangle_book_leaves.tex}; \path{lean/Taeyoung/Examples/Graph041.lean}\\
42 & \texttt{D`\{} & 01, 04, 14, 23, 24, 34 & \path{notes/pure_chordal.tex}; \path{lean/Taeyoung/Examples/Graph042.lean}\\
43 & \texttt{Dlc} & 01, 03, 04, 12, 23, 34 & \path{notes/atlas43_exact_rooted_sos.tex}; \path{lean/Taeyoung/Examples/Graph043.lean}\\
45 & \texttt{DJ\{} & 04, 12, 13, 14, 23, 24, 34 & \path{notes/clique_common_leaf_extensions.tex}; \path{lean/Taeyoung/Examples/Graph045.lean}\\
46 & \texttt{DF\{} & 03, 04, 13, 14, 23, 24, 34 & \path{notes/pure_chordal.tex}; \path{lean/Taeyoung/Examples/Graph046.lean}\\
47 & \texttt{Djs} & 01, 04, 12, 13, 14, 23, 34 & \path{notes/pure_chordal.tex}; \path{lean/Taeyoung/Examples/Graph047.lean}\\
48 & \texttt{D]w} & 02, 03, 04, 12, 13, 14, 24 & \path{lean/Taeyoung/Methods/Negative/Tensor.lean}; \path{history/paper_independent/math_paper/scripts/verify_negative_witnesses.py}; \path{lean/Taeyoung/Examples/Graph048.lean}\\
49 & \texttt{Df\{} & 01, 03, 04, 13, 14, 23, 24, 34 & \path{notes/paw_triangle_edge_cones.tex}; \path{lean/Taeyoung/Examples/Graph049.lean}\\
50 & \texttt{Dl\{} & 01, 03, 04, 12, 14, 23, 24, 34 & \path{lean/Taeyoung/Methods/Negative/Tensor.lean}; \path{history/paper_independent/math_paper/scripts/verify_negative_witnesses.py}; \path{lean/Taeyoung/Examples/Graph050.lean}\\
51 & \texttt{Dn\{} & 01, 03, 04, 12, 13, 14, 23, 24, 34 & \path{notes/pure_chordal.tex}; \path{lean/Taeyoung/Examples/Graph051.lean}\\
52 & \texttt{D\textasciitilde{}\{} & 01, 02, 03, 04, 12, 13, 14, 23, 24, 34 & \path{notes/pure_chordal.tex}; \path{lean/Taeyoung/Examples/Graph052.lean}\\
84 & \texttt{ECd\_} & 03, 04, 15, 25, 34 & \path{history/paper_independent/math_paper/main.tex}; \path{lean/Taeyoung/Examples/Graph084.lean}\\
92 & \texttt{E?Fw} & 05, 15, 25, 34, 35, 45 & \path{notes/rooted_triangle_tree_extensions.tex}; \path{lean/Taeyoung/Examples/Graph092.lean}\\
93 & \texttt{EC\{O} & 03, 04, 14, 24, 34, 35 & \path{notes/two_root_triangle_leaf_cones.tex}; \path{lean/Taeyoung/Examples/Graph093.lean}\\
94 & \texttt{E@dW} & 04, 15, 23, 34, 35, 45 & \path{notes/whiskering.tex}; \path{lean/Taeyoung/Examples/Graph094.lean}\\
95 & \texttt{EG\}?} & 04, 05, 12, 14, 24, 34 & \path{notes/mixed_rooted_triangle_branches.tex}; \path{lean/Taeyoung/Examples/Graph095.lean}\\
97 & \texttt{EYWO} & 02, 12, 13, 14, 24, 35 & \path{notes/triangle_adjacent_leaf_tail.tex}; \path{lean/Taeyoung/Examples/Graph097.lean}\\
100 & \texttt{EsCW} & 01, 02, 03, 34, 35, 45 & \path{notes/triangle_two_leaf_broom.tex}; \path{lean/Taeyoung/Examples/Graph100.lean}\\
102 & \texttt{EJe?} & 04, 05, 12, 13, 23, 34 & \path{notes/triangle_three_edge_tail.tex}; \path{lean/Taeyoung/Examples/Graph102.lean}\\
104 & \texttt{Ehd?} & 01, 04, 12, 15, 23, 34 & \path{notes/odd_cycle_one_leaf.tex}; \path{lean/Taeyoung/Examples/Graph104.lean}\\
106 & \texttt{EJaG} & 04, 05, 12, 13, 23, 45 & \path{history/paper_independent/math_paper/main.tex}; \path{lean/Taeyoung/Examples/Graph106.lean}\\
111 & \texttt{EB\{G} & 04, 13, 14, 23, 24, 34, 45 & \path{notes/forest_cone_graphon_bound.tex}; \path{lean/Taeyoung/Examples/Graph111.lean}\\
112 & \texttt{EhX\_} & 01, 12, 14, 15, 23, 24, 25 & \path{notes/two_root_triangle_leaf_cones.tex}; \path{lean/Taeyoung/Examples/Graph112.lean}\\
113 & \texttt{E\textasciicircum{}\_O} & 02, 03, 04, 12, 13, 23, 35 & \path{notes/diamond_orbit_balanced_leaves.tex}; \path{lean/Taeyoung/Examples/Graph113.lean}\\
114 & \texttt{EJwG} & 04, 12, 13, 14, 23, 24, 45 & \path{notes/page_rooted_triangle_book_leaves.tex}; \path{lean/Taeyoung/Examples/Graph114.lean}\\
115 & \texttt{E`Xg} & 01, 14, 15, 23, 24, 25, 45 & \path{notes/self_amalgam.tex}; \path{lean/Taeyoung/Examples/Graph115.lean}\\
117 & \texttt{EtaG} & 01, 02, 03, 04, 05, 23, 45 & \path{notes/forest_cone_graphon_bound.tex}; \path{lean/Taeyoung/Examples/Graph117.lean}\\
118 & \texttt{Eht?} & 01, 04, 12, 14, 15, 23, 34 & \path{notes/atlas118_compact_certificate.md}; \path{lean/Taeyoung/Examples/Graph118.lean}\\
119 & \texttt{EtoO} & 01, 02, 03, 04, 14, 23, 35 & \path{notes/bowtie_outer_leaves.tex}; \path{lean/Taeyoung/Examples/Graph119.lean}\\
120 & \texttt{EB\}?} & 04, 05, 13, 14, 23, 24, 34 & \path{notes/triangle_book_two_edge_tail.tex}; \path{lean/Taeyoung/Examples/Graph120.lean}\\
122 & \texttt{Eld?} & 01, 03, 04, 12, 15, 23, 34 & \path{notes/atlas122_compact_certificate.md}; \path{lean/Taeyoung/Examples/Graph122.lean}\\
123 & \texttt{EJy?} & 04, 05, 12, 13, 14, 23, 24 & \path{notes/triangle_book_page_two_edge_tail.tex}; \path{lean/Taeyoung/Examples/Graph123.lean}\\
124 & \texttt{Exd?} & 01, 02, 04, 12, 15, 23, 34 & \path{notes/atlas124_compact_certificate.md}; \path{lean/Taeyoung/Examples/Graph124.lean}\\
126 & \texttt{ERUO} & 02, 05, 13, 14, 23, 34, 35 & \path{notes/atlas126_compact_verification.md}; \path{notes/atlas126_triangle_c4_vertex_supporting_plane.tex}; \path{lean/Taeyoung/Examples/Graph126.lean}\\
127 & \texttt{EZEG} & 02, 05, 12, 13, 23, 34, 45 & \path{notes/atlas127_compact_certificate.md}; \path{notes/smoothed_goodman_deflagged.tex}; \path{lean/Taeyoung/Examples/Graph127.lean}\\
129 & \texttt{EheO} & 01, 04, 05, 12, 23, 34, 35 & \path{lean/Taeyoung/Methods/Negative/Tensor.lean}; \path{history/paper_independent/math_paper/scripts/verify_negative_witnesses.py}; \path{lean/Taeyoung/Examples/Graph129.lean}\\
130 & \texttt{E\{CW} & 01, 02, 03, 12, 34, 35, 45 & \path{notes/atlas130_compact_certificate.md}; \path{lean/Taeyoung/Examples/Graph130.lean}\\
133 & \texttt{E\textasciitilde{}a?} & 01, 02, 03, 04, 05, 12, 13, 23 & \path{notes/clique_common_leaf_extensions.tex}; \path{lean/Taeyoung/Examples/Graph133.lean}\\
134 & \texttt{E\textasciitilde{}\_O} & 01, 02, 03, 04, 12, 13, 23, 35 & \path{notes/clique_distributed_leaves.tex}; \path{lean/Taeyoung/Examples/Graph134.lean}\\
135 & \texttt{EzW\_} & 01, 02, 12, 13, 14, 23, 24, 25 & \path{notes/forest_cone_graphon_bound.tex}; \path{lean/Taeyoung/Examples/Graph135.lean}\\
136 & \texttt{Ejt?} & 01, 04, 12, 13, 14, 15, 23, 34 & \path{notes/forest_cone_graphon_bound.tex}; \path{lean/Taeyoung/Examples/Graph136.lean}\\
137 & \texttt{EjsG} & 01, 04, 12, 13, 14, 23, 34, 45 & \path{notes/triangle_book_page_paw_branch.tex}; \path{lean/Taeyoung/Examples/Graph137.lean}\\
138 & \texttt{Ez`\_} & 01, 02, 04, 12, 13, 15, 23, 25 & \path{notes/page_rooted_triangle_book_leaves.tex}; \path{lean/Taeyoung/Examples/Graph138.lean}\\
139 & \texttt{Eju?} & 01, 04, 05, 12, 13, 14, 23, 34 & \path{notes/triangle_book_page_paw_branch.tex}; \path{lean/Taeyoung/Examples/Graph139.lean}\\
140 & \texttt{Ev`\_} & 01, 02, 03, 04, 13, 15, 23, 25 & \path{lean/Taeyoung/Methods/Negative/Tensor.lean}; \path{history/paper_independent/math_paper/scripts/verify_negative_witnesses.py}; \path{lean/Taeyoung/Examples/Graph140.lean}\\
141 & \texttt{EXSw} & 02, 12, 14, 23, 25, 34, 35, 45 & \path{lean/Taeyoung/Methods/Negative/Tensor.lean}; \path{history/paper_independent/math_paper/scripts/verify_negative_witnesses.py}; \path{lean/Taeyoung/Examples/Graph141.lean}\\
142 & \texttt{E\textasciitilde{}AG} & 01, 02, 03, 05, 12, 13, 23, 45 & \path{notes/k4_two_edge_tail.tex}; \path{lean/Taeyoung/Examples/Graph142.lean}\\
143 & \texttt{Er`o} & 01, 02, 04, 13, 15, 23, 25, 35 & \path{lean/Taeyoung/Methods/Negative/Tensor.lean}; \path{history/paper_independent/math_paper/scripts/verify_negative_witnesses.py}; \path{lean/Taeyoung/Examples/Graph143.lean}\\
144 & \texttt{EB\}G} & 04, 05, 13, 14, 23, 24, 34, 45 & \path{notes/pure_chordal.tex}; \path{lean/Taeyoung/Examples/Graph144.lean}\\
145 & \texttt{Exe\_} & 01, 02, 04, 05, 12, 23, 25, 34 & \path{notes/c4_triangle_page_concentration.tex}; \path{lean/Taeyoung/Examples/Graph145.lean}\\
147 & \texttt{EhMg} & 01, 05, 12, 23, 24, 25, 34, 45 & \path{notes/atlas147_compact_certificate.md}; \path{lean/Taeyoung/Examples/Graph147.lean}\\
148 & \texttt{EyUG} & 01, 02, 05, 12, 13, 14, 34, 45 & \path{notes/atlas148_paw_bias_hilbert_projection.tex}; \path{lean/Taeyoung/Examples/Graph148.lean}\\
149 & \texttt{Ele\_} & 01, 03, 04, 05, 12, 23, 25, 34 & \path{lean/Taeyoung/Methods/Negative/Tensor.lean}; \path{history/paper_independent/math_paper/scripts/verify_negative_witnesses.py}; \path{lean/Taeyoung/Examples/Graph149.lean}\\
150 & \texttt{EJyG} & 04, 05, 12, 13, 14, 23, 24, 45 & \path{notes/pure_chordal.tex}; \path{lean/Taeyoung/Examples/Graph150.lean}\\
151 & \texttt{EhdW} & 01, 04, 12, 15, 23, 34, 35, 45 & \path{notes/atlas151_compact_certificate.md}; \path{lean/Taeyoung/Examples/Graph151.lean}\\
152 & \texttt{EhNG} & 01, 05, 12, 15, 23, 24, 34, 45 & \path{notes/atlas_152_fractional_fibre_local_counterexample.tex}; \path{lean/Taeyoung/Examples/Graph152.lean}\\
153 & \texttt{Ehf\_} & 01, 04, 05, 12, 15, 23, 25, 34 & \path{notes/atlas153_compact_certificate.md}; \path{lean/Taeyoung/Examples/Graph153.lean}\\
156 & \texttt{E\textasciitilde{}H\_} & 01, 02, 03, 12, 13, 15, 23, 24, 25 & \path{notes/paw_triangle_edge_cones.tex}; \path{lean/Taeyoung/Examples/Graph156.lean}\\
157 & \texttt{E\textasciitilde{}`\_} & 01, 02, 03, 04, 12, 13, 15, 23, 25 & \path{notes/atlas157_compact_certificate.md}; \path{lean/Taeyoung/Examples/Graph157.lean}\\
158 & \texttt{El\{G} & 01, 03, 04, 12, 14, 23, 24, 34, 45 & \path{lean/Taeyoung/Methods/Negative/Tensor.lean}; \path{history/paper_independent/math_paper/scripts/verify_negative_witnesses.py}; \path{lean/Taeyoung/Examples/Graph158.lean}\\
159 & \texttt{EZSw} & 02, 12, 13, 14, 23, 25, 34, 35, 45 & \path{lean/Taeyoung/Methods/Negative/Tensor.lean}; \path{history/paper_independent/math_paper/scripts/verify_negative_witnesses.py}; \path{lean/Taeyoung/Examples/Graph159.lean}\\
160 & \texttt{E\textasciitilde{}@g} & 01, 02, 03, 12, 13, 15, 23, 25, 45 & \path{notes/atlas160_k4_paw_edge_supporting_plane.tex}; \path{lean/Taeyoung/Examples/Graph160.lean}\\
161 & \texttt{E?\textasciitilde{}w} & 04, 05, 14, 15, 24, 25, 34, 35, 45 & \path{notes/pure_chordal.tex}; \path{lean/Taeyoung/Examples/Graph161.lean}\\
162 & \texttt{E|e\_} & 01, 02, 03, 04, 05, 12, 23, 25, 34 & \path{notes/pure_chordal.tex}; \path{lean/Taeyoung/Examples/Graph162.lean}\\
163 & \texttt{EyuG} & 01, 02, 04, 05, 12, 13, 14, 34, 45 & \path{notes/pure_chordal.tex}; \path{lean/Taeyoung/Examples/Graph163.lean}\\
164 & \texttt{EyVG} & 01, 02, 05, 12, 13, 14, 15, 34, 45 & \path{notes/pure_chordal.tex}; \path{lean/Taeyoung/Examples/Graph164.lean}\\
165 & \texttt{E\textasciitilde{}aG} & 01, 02, 03, 04, 05, 12, 13, 23, 45 & \path{notes/paw_triangle_edge_cones.tex}; \path{lean/Taeyoung/Examples/Graph165.lean}\\
166 & \texttt{ElfO} & 01, 03, 04, 05, 12, 15, 23, 34, 35 & \path{notes/turan_local_and_high_density_negative_tests.tex}; \path{lean/Taeyoung/Examples/Graph166.lean}\\
167 & \texttt{E\textasciicircum{}eG} & 02, 03, 04, 05, 12, 13, 23, 34, 45 & \path{notes/pure_chordal.tex}; \path{lean/Taeyoung/Examples/Graph167.lean}\\
168 & \texttt{E\textasciicircum{}MG} & 02, 03, 05, 12, 13, 23, 24, 34, 45 & \path{notes/atlas168_compact_certificate.md}; \path{lean/Taeyoung/Examples/Graph168.lean}\\
169 & \texttt{Exf\_} & 01, 02, 04, 05, 12, 15, 23, 25, 34 & \path{notes/atlas169_compact_certificate.md}; \path{lean/Taeyoung/Examples/Graph169.lean}\\
170 & \texttt{EO\textasciitilde{}o} & 02, 04, 05, 14, 15, 24, 25, 34, 35 & \path{lean/Taeyoung/Methods/Negative/Tensor.lean}; \path{history/paper_independent/math_paper/scripts/verify_negative_witnesses.py}; \path{lean/Taeyoung/Examples/Graph170.lean}\\
171 & \texttt{Ehew} & 01, 04, 05, 12, 23, 25, 34, 35, 45 & \path{notes/atlas171_compact_certificate.md}; \path{lean/Taeyoung/Examples/Graph171.lean}\\
172 & \texttt{Elf\_} & 01, 03, 04, 05, 12, 15, 23, 25, 34 & \path{notes/turan_local_and_high_density_negative_tests.tex}; \path{lean/Taeyoung/Examples/Graph172.lean}\\
173 & \texttt{ElMg} & 01, 03, 05, 12, 23, 24, 25, 34, 45 & \path{lean/Taeyoung/Methods/Negative/Tensor.lean}; \path{history/paper_independent/math_paper/scripts/verify_negative_witnesses.py}; \path{lean/Taeyoung/Examples/Graph173.lean}\\
174 & \texttt{EtTg} & 01, 02, 03, 14, 15, 23, 25, 34, 45 & \path{notes/atlas174_compact_certificate.md}; \path{lean/Taeyoung/Examples/Graph174.lean}\\
177 & \texttt{En\{G} & 01, 03, 04, 12, 13, 14, 23, 24, 34, 45 & \path{notes/six_verified_base_cones.tex}; \path{lean/Taeyoung/Examples/Graph177.lean}\\
178 & \texttt{En\}?} & 01, 03, 04, 05, 12, 13, 14, 23, 24, 34 & \path{notes/atlas178_half_degree_weighted_k4.tex}; \path{lean/Taeyoung/Examples/Graph178.lean}\\
179 & \texttt{E\_\textasciitilde{}w} & 01, 04, 05, 14, 15, 24, 25, 34, 35, 45 & \path{notes/six_verified_base_cones.tex}; \path{lean/Taeyoung/Examples/Graph179.lean}\\
180 & \texttt{EjtW} & 01, 04, 12, 13, 14, 15, 23, 34, 35, 45 & \path{notes/two_root_triangle_leaf_cones.tex}; \path{lean/Taeyoung/Examples/Graph180.lean}\\
181 & \texttt{E\textasciicircum{}mG} & 02, 03, 04, 05, 12, 13, 23, 24, 34, 45 & \path{notes/atlas181_compact_certificate.md}; \path{lean/Taeyoung/Examples/Graph181.lean}\\
182 & \texttt{E\textasciicircum{}Mg} & 02, 03, 05, 12, 13, 23, 24, 25, 34, 45 & \path{lean/Taeyoung/Methods/Negative/Tensor.lean}; \path{history/paper_independent/math_paper/scripts/verify_negative_witnesses.py}; \path{lean/Taeyoung/Examples/Graph182.lean}\\
183 & \texttt{EjvG} & 01, 04, 05, 12, 13, 14, 15, 23, 34, 45 & \path{notes/six_verified_base_cones.tex}; \path{lean/Taeyoung/Examples/Graph183.lean}\\
184 & \texttt{Elfo} & 01, 03, 04, 05, 12, 15, 23, 25, 34, 35 & \path{lean/Taeyoung/Methods/Negative/Tensor.lean}; \path{history/paper_independent/math_paper/scripts/verify_negative_witnesses.py}; \path{lean/Taeyoung/Examples/Graph184.lean}\\
185 & \texttt{Exv\_} & 01, 02, 04, 05, 12, 14, 15, 23, 25, 34 & \path{notes/atlas185_compact_certificate.md}; \path{lean/Taeyoung/Examples/Graph185.lean}\\
186 & \texttt{ErXw} & 01, 02, 13, 14, 15, 23, 24, 25, 35, 45 & \path{lean/Taeyoung/Methods/Negative/Tensor.lean}; \path{history/paper_independent/math_paper/scripts/verify_negative_witnesses.py}; \path{lean/Taeyoung/Examples/Graph186.lean}\\
187 & \texttt{Ehfw} & 01, 04, 05, 12, 15, 23, 25, 34, 35, 45 & \path{notes/clique_oddcycle_join.tex}; \path{lean/Taeyoung/Examples/Graph187.lean}\\
188 & \texttt{EzNG} & 01, 02, 05, 12, 13, 15, 23, 24, 34, 45 & \path{notes/atlas188_compact_certificate.md}; \path{lean/Taeyoung/Examples/Graph188.lean}\\
189 & \texttt{E\textasciicircum{}NG} & 02, 03, 05, 12, 13, 15, 23, 24, 34, 45 & \path{lean/Taeyoung/Methods/Negative/Tensor.lean}; \path{history/paper_independent/math_paper/scripts/verify_negative_witnesses.py}; \path{lean/Taeyoung/Examples/Graph189.lean}\\
190 & \texttt{EyUw} & 01, 02, 05, 12, 13, 14, 25, 34, 35, 45 & \path{lean/Taeyoung/Methods/Negative/Tensor.lean}; \path{history/paper_independent/math_paper/scripts/verify_negative_witnesses.py}; \path{lean/Taeyoung/Examples/Graph190.lean}\\
191 & \texttt{E\textasciitilde{}|?} & 01, 02, 03, 04, 12, 13, 14, 15, 23, 24, 34 & \path{notes/clique_common_leaf_extensions.tex}; \path{lean/Taeyoung/Examples/Graph191.lean}\\
192 & \texttt{E\textasciitilde{}Xo} & 01, 02, 03, 12, 13, 14, 15, 23, 24, 25, 35 & \path{notes/six_verified_base_cones.tex}; \path{lean/Taeyoung/Examples/Graph192.lean}\\
193 & \texttt{En\}G} & 01, 03, 04, 05, 12, 13, 14, 23, 24, 34, 45 & \path{notes/page_rooted_triangle_book_leaves.tex}; \path{lean/Taeyoung/Examples/Graph193.lean}\\
194 & \texttt{E\textasciitilde{}wW} & 01, 02, 03, 04, 12, 13, 14, 23, 24, 35, 45 & \path{notes/atlas194_compact_certificate.md}; \path{lean/Taeyoung/Examples/Graph194.lean}\\
195 & \texttt{EyVw} & 01, 02, 05, 12, 13, 14, 15, 25, 34, 35, 45 & \path{notes/pure_chordal.tex}; \path{lean/Taeyoung/Examples/Graph195.lean}\\
196 & \texttt{ER\textasciitilde{}g} & 02, 04, 05, 13, 14, 15, 23, 24, 25, 34, 45 & \path{notes/open_conditional_implications.tex}; \path{lean/Taeyoung/Examples/Graph196.lean}\\
197 & \texttt{E\}\textasciicircum{}G} & 01, 02, 03, 05, 12, 13, 14, 15, 24, 34, 45 & \path{lean/Taeyoung/Methods/Negative/Tensor.lean}; \path{history/paper_independent/math_paper/scripts/verify_negative_witnesses.py}; \path{lean/Taeyoung/Examples/Graph197.lean}\\
198 & \texttt{Ep\textasciitilde{}o} & 01, 02, 04, 05, 14, 15, 23, 24, 25, 34, 35 & \path{lean/Taeyoung/Methods/Negative/Tensor.lean}; \path{history/paper_independent/math_paper/scripts/verify_negative_witnesses.py}; \path{lean/Taeyoung/Examples/Graph198.lean}\\
199 & \texttt{El\textasciicircum{}g} & 01, 03, 05, 12, 14, 15, 23, 24, 25, 34, 45 & \path{notes/atlas199_compact_certificate.md}; \path{lean/Taeyoung/Examples/Graph199.lean}\\
200 & \texttt{E\textasciitilde{}\{W} & 01, 02, 03, 04, 12, 13, 14, 23, 24, 34, 35, 45 & \path{notes/six_verified_base_cones.tex}; \path{lean/Taeyoung/Examples/Graph200.lean}\\
201 & \texttt{E\textasciitilde{}z\_} & 01, 02, 03, 04, 05, 12, 13, 14, 15, 23, 24, 25 & \path{notes/pure_chordal.tex}; \path{lean/Taeyoung/Examples/Graph201.lean}\\
202 & \texttt{Ep\textasciitilde{}w} & 01, 02, 04, 05, 14, 15, 23, 24, 25, 34, 35, 45 & \path{notes/pure_chordal.tex}; \path{lean/Taeyoung/Examples/Graph202.lean}\\
203 & \texttt{E\textasciitilde{}\textasciicircum{}G} & 01, 02, 03, 05, 12, 13, 14, 15, 23, 24, 34, 45 & \path{notes/atlas203_compact_certificate.md}; \path{lean/Taeyoung/Examples/Graph203.lean}\\
204 & \texttt{EznW} & 01, 02, 04, 05, 12, 13, 15, 23, 24, 34, 35, 45 & \path{lean/Taeyoung/Methods/Negative/Tensor.lean}; \path{history/paper_independent/math_paper/scripts/verify_negative_witnesses.py}; \path{lean/Taeyoung/Examples/Graph204.lean}\\
205 & \texttt{E\textasciitilde{}\textasciitilde{}G} & 01, 02, 03, 04, 05, 12, 13, 14, 15, 23, 24, 34, 45 & \path{notes/six_verified_base_cones.tex}; \path{lean/Taeyoung/Examples/Graph205.lean}\\
206 & \texttt{E\textasciitilde{}nW} & 01, 02, 03, 04, 05, 12, 13, 15, 23, 24, 34, 35, 45 & \path{notes/turan_local_and_high_density_negative_tests.tex}; \path{lean/Taeyoung/Examples/Graph206.lean}\\
207 & \texttt{Ez\textasciitilde{}w} & 01, 02, 04, 05, 12, 13, 14, 15, 23, 24, 25, 34, 35, 45 & \path{notes/pure_chordal.tex}; \path{lean/Taeyoung/Examples/Graph207.lean}\\
208 & \texttt{E\textasciitilde{}\textasciitilde{}w} & 01, 02, 03, 04, 05, 12, 13, 14, 15, 23, 24, 25, 34, 35, 45 & \path{notes/pure_chordal.tex}; \path{lean/Taeyoung/Examples/Graph208.lean}\\
\end{longtable}
\end{landscape}


\section{Exact negative evaluations}
\label{app:negative}
\begin{landscape}
\scriptsize
\begin{longtable}{r p{4.0cm} l p{4.1cm} p{4.1cm} p{4.8cm}}
\caption{Exact counterexamples for the 23 negative rows.}\label{tab:negative-witnesses}\\
\toprule Atlas & witness & $p$ & $t(H,W)$ & $\Phi_H(p)$ & $t(H,W)-\Phi_H(p)$\\ \midrule
\endfirsthead
\toprule Atlas & witness & $p$ & $t(H,W)$ & $\Phi_H(p)$ & $t(H,W)-\Phi_H(p)$\\ \midrule
\endhead
\bottomrule\endfoot
48 & $T_{4}\otimes T_{4}$ & $\frac{9}{16}$ & $\frac{225}{16384}$ & $\frac{477}{32768}$ & $\frac{-27}{32768}$\\
50 & $T_{3}\otimes T_{5}$ & $\frac{8}{15}$ & $\frac{56}{16875}$ & $\frac{344}{50625}$ & $\frac{-176}{50625}$\\
129 & $T_{6}\otimes T_{8}$ & $\frac{35}{48}$ & $\frac{1171205}{10616832}$ & $\frac{14084455}{127401984}$ & $\frac{-29995}{127401984}$\\
140 & $T_{4}\otimes T_{4}$ & $\frac{9}{16}$ & $\frac{2025}{262144}$ & $\frac{4293}{524288}$ & $\frac{-243}{524288}$\\
141 & $T_{4}\otimes T_{4}$ & $\frac{9}{16}$ & $\frac{2025}{262144}$ & $\frac{4293}{524288}$ & $\frac{-243}{524288}$\\
143 & $T_{4}\otimes T_{4}$ & $\frac{9}{16}$ & $\frac{2025}{262144}$ & $\frac{4293}{524288}$ & $\frac{-243}{524288}$\\
149 & $T_{4}\otimes T_{5}$ & $\frac{3}{5}$ & $\frac{6579}{400000}$ & $\frac{57}{3125}$ & $\frac{-717}{400000}$\\
152 & five-step fractional-fibre perturbation & $\frac{18750093281}{37500000000}$ & $\frac{185159623403}{60000000000000000000000}$ & \resizebox{3.95cm}{!}{$\frac{19119261293088692808932739896407207170401}{6179809570312500000000000000000000000000000000000000}$} & \resizebox{4.65cm}{!}{$\frac{-48407747400063658542114896407207170401}{6179809570312500000000000000000000000000000000000000}$}\\
158 & $T_{3}\otimes T_{5}$ & $\frac{8}{15}$ & $\frac{448}{253125}$ & $\frac{2752}{759375}$ & $\frac{-1408}{759375}$\\
159 & $T_{3}\otimes T_{5}$ & $\frac{8}{15}$ & $\frac{448}{253125}$ & $\frac{2752}{759375}$ & $\frac{-1408}{759375}$\\
166 & four-step two-scale perturbation & $\frac{9}{16}$ & $\frac{465}{262144}$ & $\frac{477}{262144}$ & $\frac{-3}{65536}$\\
170 & $T_{3}\otimes T_{5}$ & $\frac{8}{15}$ & $\frac{1088}{253125}$ & $\frac{3536}{759375}$ & $\frac{-272}{759375}$\\
172 & four-step two-scale perturbation & $\frac{9}{16}$ & $\frac{3}{2048}$ & $\frac{477}{262144}$ & $\frac{-93}{262144}$\\
173 & $T_{4}\otimes T_{6}$ & $\frac{5}{8}$ & $\frac{5}{384}$ & $\frac{55}{4096}$ & $\frac{-5}{12288}$\\
182 & $T_{4}\otimes T_{4}$ & $\frac{9}{16}$ & $\frac{81}{65536}$ & $\frac{351}{262144}$ & $\frac{-27}{262144}$\\
184 & $T_{4}\otimes T_{4}$ & $\frac{9}{16}$ & $\frac{81}{65536}$ & $\frac{351}{262144}$ & $\frac{-27}{262144}$\\
186 & $T_{4}\otimes T_{4}$ & $\frac{9}{16}$ & $\frac{9}{4096}$ & $\frac{99}{32768}$ & $\frac{-27}{32768}$\\
189 & $T_{5}\otimes T_{5}$ & $\frac{16}{25}$ & $\frac{97344}{9765625}$ & $\frac{97888}{9765625}$ & $\frac{-544}{9765625}$\\
190 & $T_{4}\otimes T_{4}$ & $\frac{9}{16}$ & $\frac{9}{4096}$ & $\frac{99}{32768}$ & $\frac{-27}{32768}$\\
197 & $T_{4}\otimes T_{6}$ & $\frac{5}{8}$ & $\frac{1075}{331776}$ & $\frac{55}{16384}$ & $\frac{-155}{1327104}$\\
198 & $T_{4}\otimes T_{6}$ & $\frac{5}{8}$ & $\frac{1075}{331776}$ & $\frac{55}{16384}$ & $\frac{-155}{1327104}$\\
204 & $T_{4}\otimes T_{6}$ & $\frac{5}{8}$ & $\frac{85}{41472}$ & $\frac{5}{2048}$ & $\frac{-65}{165888}$\\
206 & three-step one-diagonal perturbation & $\frac{25}{36}$ & $\frac{77731841}{48922361856}$ & $\frac{16975}{10077696}$ & $\frac{-126187493}{1320903770112}$\\
\end{longtable}
\end{landscape}


\section{Retained Lean build measurements}
\label{app:verification}
\begin{table}[p]
\centering
\small
\caption{Measured sequential Lean builds for the twenty compact-certificate rows.  Time includes the fresh row build and the installed-source check when reported.}
\label{tab:verification-cost}
\begin{tabular}{rrrr}
\toprule Atlas & source files & time (min) & peak memory (GiB)\\ \midrule
118 & 89 & 25.7 & 9.09\\
122 & 89 & 25.9 & 9.09\\
124 & 89 & 25.8 & 9.09\\
126 & 100 & 45.3 & 9.73\\
127 & 97 & 43.6 & 11.90\\
130 & 100 & 34.6 & 9.21\\
147 & 97 & 43.8 & 11.84\\
151 & 97 & 46.2 & 11.83\\
153 & 98 & 40.9 & 11.90\\
157 & 89 & 27.0 & 9.20\\
168 & 97 & 44.7 & 11.84\\
169 & 89 & 26.9 & 9.23\\
171 & 98 & 42.4 & 11.88\\
174 & 98 & 41.0 & 11.91\\
181 & 89 & 26.9 & 9.23\\
185 & 89 & 28.6 & 9.23\\
188 & 98 & 41.0 & 11.90\\
194 & 89 & 26.8 & 9.21\\
199 & 89 & 26.8 & 9.23\\
203 & 100 & 37.5 & 9.26\\
\midrule
Total & -- & 701.5 & --\\
\bottomrule
\end{tabular}
\end{table}


\bibliographystyle{plain}
\bibliography{references}

\end{document}
